\Needspace{5\baselineskip}\textbf{Problems}

\begin{enumerate}[label=\thesection.\arabic*.,ref=\thesection.\arabic*,start=5,series=book-problems]

\item\label{ex:mean-zero-neumann-variational-compatibility} Assume that $\Omega$ is bounded, connected, and has the $1$-extension property. Let\[V=\left\{u\in H^1(\Omega)\ \middle|\ \int_\Omega u(\bm{x})\,\mathrm{d}\bm{x}=0\right\}.\](a) Show that for each $f\in L^2(\Omega)$ there is a unique $u\in V$ such that
\begin{equation}\label{eq:mean-zero-neumann-variational-problem}
\int_\Omega\nabla u\cdot\nabla v=\int_\Omega fv\quad\text{for all }v\in V.
\end{equation}
(See Problem~\ref{ex:mean-zero-sobolev-poincare-inequality}.)

(b) Explain why it is appropriate to regard \eqref{eq:mean-zero-neumann-variational-problem} as a weak form of the Neumann problem
\begin{equation}\label{eq:mean-zero-neumann-strong-problem}
\begin{aligned}
-\Delta u&=f&&\text{in }\Omega\\
\frac{\partial u}{\partial n}&=0&&\text{on }\partial\Omega
\end{aligned}
\end{equation}
if $\int_\Omega f=0$.

(c) If $\int_\Omega f\ne0$, is it still reasonable to call the solution of \eqref{eq:mean-zero-neumann-variational-problem} a solution of \eqref{eq:mean-zero-neumann-strong-problem}? Explain.

\item\label{ex:scalar-dirichlet-complementing-condition} Show that Dirichlet boundary conditions for scalar elliptic PDEs always satisfy the complementing condition.

\item\label{ex:dirichlet-system-triangular-reduction-details} Fill in the details for the proof of Lemma~\ref{lem:dirichlet-system-tangential-triangular-reduction}.

\item\label{ex:agmon-condition-implies-complementing-natural-boundaries} Suppose that Agmon\textquotesingle s condition holds. Show that the complementing condition is satisfied for \eqref{eq:agmon-essential-natural-boundary-value-problem}. Hint: Apply \eqref{eq:agmon-constrained-form-green-boundary-identity} on a half-space.

\item\label{ex:inhomogeneous-essential-natural-boundaries-weak-form} Formulate a weak form of \eqref{eq:agmon-essential-natural-boundary-value-problem} when the boundary conditions are allowed to be inhomogeneous.

\item\label{ex:stokes-traction-complementing-condition} Show that the ``traction boundary conditions\textquotesingle\textquotesingle\ $(\nabla\bm{u}+(\nabla\bm{u})^{\mathrm{T}})\cdot\bm{n}-p\bm{n}=0$ satisfy the complementing condition for the Stokes system.

\item\label{ex:scalar-dirichlet-zero-fredholm-index} Show that a scalar elliptic operator with Dirichlet conditions has Fredholm index $0$. Hint: Show that the adjoint problem also has Dirichlet conditions.

\end{enumerate}
